\section{The Architecture of Transition}
\label{sec:architecture_of_transition}

Any persistent system partitioned into nested subsystems must eventually transfer information across the boundaries between them. In \textit{Informational Energetics: Nested Persistence}~\cite{meyer_ienested_2026}, it was established that when internal complexity exceeds the capacity of a single control loop, a system must structurally partition itself. However, partitioning solves only the problem of structural containment. For the partitioned system to function as a coherent whole, information must flow dynamically across the interfaces between sectors.

\subsection{The Dissipation Density of Boundary Crossing}

When information crosses from partition $i$ to partition $j$, the transition incurs an irreducible entropic cost. As established in \textit{Informational Energetics: Entropic Action}~\cite{meyer_ieaction_2026}, the instantaneous dissipation density is: 
$$ \mathcal{D}_{ij} = \xi \cdot \mathbf{K}_{i\to j} \cdot f $$ 
where:
\begin{itemize}
    \item $\xi$ is the \textbf{transition impedance}: the fundamental resistance of the substrate to state transfer (governed by the intrinsic noise floor and physical constants of the medium).
    \item $\mathbf{K}_{i\to j}$ is the \textbf{topological mismatch}: the geometric stress required to deform the information payload from the source partition's capacity constraints to the destination's coordinate system.
    \item $f$ is the \textbf{transition flux}: the probability density of the transfer event per unit time.
\end{itemize}

The multiplicative form follows from the statistical independence of channel friction, payload complexity, and transmission rate~\cite{meyer_ieaction_2026}. By the Principle of Least Entropic Action, persistent transitions are those trajectories that minimize the path integral of $\mathcal{D}_{ij}$.

\subsection{The Three Transition Archetypes}

 (continuous transfer across adjacent boundaries, limited by interface friction), 
 (hierarchical expansion of structural scope, engaging higher-order capacity;)  (discontinuous transfer across non-adjacent tiers, bypassing intermediate structure).

\begin{enumerate}
    \item \textbf{Adjacent Transfer (Type I):} Adjacent, continuous transitions governed by localized interface friction. The information payload deforms only at the shared boundary between adjacent partitions, incurring an entropic cost proportional to the localized surface area of the interface.
    
    \item \textbf{Escalation (Type II):} Volumetric promotions engaging the full field impedance. The payload expands from a lower-dimensional to a higher-dimensional partition within the nested hierarchy, performing work to scale its operational capacity against the macroscopic substrate resistance.
    
    \item \textbf{Direct Routing (Type III):} Non-adjacent transitions across a macroscopic geometric barrier. The payload bypasses intermediate partitions via a discrete discontinuous step, violating continuous deformation and incurring maximal localized impedance to map directly to a remote coordinate space.
\end{enumerate}

These three archetypes exhaust the stable transition geometry on a finite-capacity, discrete substrate. A fourth logical possibility—the complete dissolution of the partition boundary—triggers \textit{Causal Aliasing}, wherein the system loses structural identity. This is not a stable transition but a catastrophic phase collapse. While macroscopic persistent systems interacting with external environments may exhibit localized `Structural Shedding' as a survival heuristic, the analysis of such boundary destruction is deferred to future work.